\section{Nhìn mạng như một hệ động lực}
\label{sec:ch03-he-dong-luc}

\index{hệ động lực!miền hút}%
Hai mục vừa rồi trả lời câu hỏi \enquote{khi nào} bằng đại số tuyến tính. Bài
báo còn trả lời nó lần thứ hai, bằng hình học, và lần này câu trả lời giải
thích được một chuyện mà đại số tuyến tính không giải thích: vì sao huấn luyện
hỏng đột ngột thay vì hỏng dần.

Coi mạng không có input là một \tn{hệ động lực}{dynamical system}. Với mỗi bộ
tham số \(\theta\), trạng thái chạy theo thời gian và tiến về một
\tn{điểm hút}{attractor}. Không gian trạng thái chia thành các
\tn{miền hút}{basin of attraction}, mỗi miền một điểm hút. Khi \(\theta\) đổi
chậm, hành vi tiệm cận cũng đổi chậm, gần như ở mọi nơi. Gần như, chứ không
phải mọi nơi: có những chỗ mà một thay đổi nhỏ của \(\theta\) làm điểm hút xuất
hiện, biến mất hoặc đổi hình. Những chỗ ấy là
\tn{điểm rẽ nhánh}{bifurcation}.

Bài dựa cách nhìn này trên công trình của Doya từ 1993, và mở rộng nó theo hai
hướng. Thứ nhất, họ phân biệt hai sự kiện: vượt qua ranh giới giữa hai miền hút,
và vượt qua một điểm rẽ nhánh. Chỉ sự kiện thứ nhất là \tn{điều kiện đủ}{sufficient condition} để \tn{độ dốc}{gradient}
nổ. Bài nói thẳng rằng vượt điểm rẽ nhánh \emph{không} phải \tn{điều kiện cần}{necessary condition} cũng
không phải điều kiện đủ, vì điểm rẽ nhánh là sự kiện toàn cục và có thể chẳng
ảnh hưởng gì tại chỗ mô hình đang đứng. Thứ hai, họ xử lý được cả mô hình có
input, thứ mà lý thuyết hệ động lực cổ điển không nhận, bằng cách tách ánh xạ
thành một phần cố định và một phần thay đổi theo thời gian.

\subsection*{Vách trong mặt lỗi}

\index{mặt lỗi!vách}%
Phần đáng giá nhất của mục này là chỗ bài hạ bài toán xuống hai tham số để vẽ
được \tn{mặt lỗi}{error surface}. Một đơn vị ẩn, không input:

\begin{equation}
  a_t = w\, \sigma(a_{t-1}) + b,
  \label{eq:ch03-mo-hinh-mot-don-vi}
\end{equation}

với \(\sigma\) là sigmoid, \(a_0 = 0.5\), và chi phí đọc đúng một lần ở bước 50:
\(L_{50} = (\sigma(a_{50}) - 0.7)^2\). Toàn bộ mô hình là hai số,
\(w\) và \(b\), nên mặt lỗi vẽ ra được.

Bài vẽ nó và chỉ vào một bức tường gần như dựng đứng. Hình vẽ thì thuyết phục,
nhưng hình vẽ không phải số đo, và cái tôi muốn biết là bức tường ấy hẹp đến
đâu. Tôi quét \(b\) ngang qua nó tại \(w = 5.0\), bước 0.001:

\begin{minted}[fontsize=\footnotesize]{text}
         b          cost        ||grad||
   -2.6140    0.34182245    2.457457e-01
   -2.6130    0.34157104    2.942246e-01
   -2.6120    0.00847633    9.281034e+00
   -2.6110    0.01020121    5.137097e-01
   -2.6100    0.01049210    2.777983e-01
\end{minted}

Giữa hai điểm cách nhau 0.001, chi phí rơi 0.33309471. Hình~\ref{fig:ch03-vach-mat-loi}
vẽ lát cắt đó.

\begin{figure}[htbp]
  \centering
  % Mặt lỗi của mô hình một đơn vị ẩn, cắt dọc theo b tại w = 5.0.
% Số liệu: research/2026-08-ch03-phan-tich.md, mục "Đo 3".
% Toạ độ: x_cm = (b + 2.620) * 500, y_cm = cost * 13.889.
% Mono-safe: chỉ dùng nét liền, nét đứt và mức xám.
\begin{tikzpicture}[
  duongcong/.style = {line width = 0.9pt, black},
  truc/.style      = {->, >={Stealth[length=2mm]}, gray!70, line width = 0.5pt},
  ghichu/.style    = {font = \footnotesize},
  moc/.style       = {densely dashed, gray!60, line width = 0.5pt},
]

  % Trục
  \draw[truc] (-0.4, 0) -- (10.6, 0) node[right, ghichu, black] {\(b\)};
  \draw[truc] (0, -0.4) -- (0, 5.4) node[above, ghichu, black] {\(L_{50}\)};

  % Vạch trục x
  \foreach \bx/\blabel in {0/-2.620, 5/-2.610, 10/-2.600} {
    \draw[gray!70] (\bx, 0) -- (\bx, -0.12);
    \node[below, ghichu] at (\bx, -0.12) {\(\blabel\)};
  }
  % Vạch trục y
  \foreach \by/\blabel in {0.139/0.01, 4.861/0.35} {
    \draw[gray!70] (0, \by) -- (-0.12, \by);
    \node[left, ghichu] at (-0.12, \by) {\(\blabel\)};
  }

  % Nhánh trái: gần như phẳng
  \draw[duongcong] (0, 4.768) -- (1.5, 4.758) -- (3.0, 4.748) -- (3.5, 4.744);

  % Vách
  \draw[duongcong] (3.5, 4.744) -- (4.0, 0.118);

  % Nhánh phải: phẳng, dốc lên rất nhẹ
  \draw[duongcong] (4.0, 0.118) -- (4.5, 0.142) -- (5.0, 0.146) -- (7.5, 0.158)
                                -- (10.0, 0.168);

  % Điểm dốc nhất. x = (-2.61234136 + 2.620) * 500 = 3.8293.
  \draw[moc] (3.829, 0) -- (3.829, 5.0);
  \node[ghichu, align = left, anchor = west] at (4.3, 4.6)
    {đỉnh dốc tại \(b = -2.6123413600\)\\ \(\norm{\nabla L_{50}} = 6.773125 \times 10^{3}\)};

  % Vùng phẳng nằm ngoài lát cắt này, ở b = -2.0, nên chú thích không có mốc
  % chỉ vào đâu cả: vẽ một mốc ở đây là chỉ vào một điểm mà con số không lấy từ đó.
  \node[ghichu, anchor = north east, align = right] at (10.5, 2.9)
    {tại \(b = -2.0\), ngoài lát cắt này:\\
     \(\norm{\nabla L_{50}} \approx 5.549566 \times 10^{-2}\)};

  % Nhãn hai nhánh
  \node[ghichu, anchor = south west] at (0.15, 4.85) {\(L_{50} \approx 0.342\)};
  \node[ghichu, anchor = south west] at (4.6, 0.32) {\(L_{50} \approx 0.012\)};

\end{tikzpicture}

  \caption{Lát cắt mặt lỗi dọc theo \(b\) tại \(w = 5.0\). Hai nhánh gần như
    phẳng, nối nhau bằng một vách rộng chưa tới một phần nghìn của \(b\). Độ dốc
    tại đỉnh vách lớn hơn độ dốc vùng phẳng năm bậc độ lớn, nên một bước
    gradient descent dùng chung một learning rate cho cả hai vùng không thể vừa
    an toàn ở đây vừa tiến được ở kia.}
  \label{fig:ch03-vach-mat-loi}
\end{figure}

Lưới 0.001 vẫn còn quá thô. Tìm bằng lưới lồng nhau bốn vòng thì đỉnh thật nằm
ở \(b = -2.6123413600\), và ở đó chuẩn gradient là 6.773125 nhân mười mũ ba.

\begin{measured}{độ dốc tại vách so với vùng phẳng}
\begin{minted}[fontsize=\footnotesize]{text}
    w          b at peak    max ||grad||    cost there
  4.6      -2.3642758000    3.157208e+02    0.13157964
  4.8      -2.4880230200    1.396279e+03    0.14229738
  5.0      -2.6123413600    6.773125e+03    0.14960413
  5.2      -2.7367973600    3.424156e+04    0.15545694
  5.4      -2.8612821400    1.736738e+05    0.16071180
\end{minted}

Tại \(w = 5.0\), chuẩn gradient ở vùng phẳng \(b = -2.0\) là
5.549566 nhân mười mũ trừ hai. Tỷ số giữa đỉnh vách và vùng phẳng là
1.2205 nhân mười mũ năm, tức 5.0865 bậc độ lớn.

Đỉnh vách di chuyển gần như tuyến tính theo \(w\), khoảng \(-0.62\) đơn vị
\(b\) trên mỗi đơn vị \(w\), và độ dốc tại đỉnh tăng khoảng năm lần sau mỗi
0.2 của \(w\).

Đo bằng \code{experiments/ch03\_surface.py} tại tag \code{ch03}.
\end{measured}

Hai chi tiết trong lúc đo đáng nói, vì cả hai đều là cách dễ dàng để đo sai.

Lần quét đầu tiên tôi dùng bước 0.01 và không thấy vách đâu cả: chuẩn gradient
lớn nhất trên cả cửa sổ chỉ 0.2777983, và mặt lỗi trông như một cái dốc thoải.
Lưới nhảy qua chỗ dốc. Con số ấy thấp hơn đỉnh thật bốn bậc độ lớn, nên nó không
phải một ước lượng thô của đỉnh mà là một con số khác hẳn.

Chi tiết thứ hai: tại \(w = 5.4\), đỉnh nằm ở \(b = -2.8612821400\), tức
\emph{ngoài} cửa sổ mà chính bài báo vẽ. Giữ nguyên cửa sổ của bài thì tìm kiếm
trả về 0.1355472 tại \(b = -2.80816320\), sát mép trái, và con số đó vô nghĩa.

\subsection*{Vì sao vách làm hỏng một bước gradient descent}

\index{độ cong!tại vách}%
Giả thuyết của bài là chỗ này: khi gradient nổ thì \tn{độ cong}{curvature} cũng
nổ theo, nên hình dạng cục bộ ở đó là một thung lũng có một bức tường dựng, chứ
chưa bao giờ là một cái dốc đều. Nếu hướng gradient và hướng cong lớn nhất trùng nhau thì mặt lỗi
dốc đứng theo phương vuông góc với tường, và một bước gradient descent chạm
tường sẽ nhảy ngang qua thung lũng, ra khỏi vùng nó đang ở.

Điều làm bước nhảy ấy tai hại không phải hướng mà là độ dài. Hướng vẫn là hướng
giảm. Chỉ là nó dài tới mức thông tin cục bộ dùng để chọn nó đã hết hiệu lực từ
lâu trước khi bước đi kết thúc.

Nói cách đó thì cách sửa gần như tự viết ra: giữ hướng, chặn độ dài.
